\section{Related work}\label{sec:related-work}

We situate the contribution within process modelling broadly rather than against categorical semantics alone~\cite{chechikFormalMethodsScope2025}. The threads on notation diversity, BPMN semantics, and object-centric mining carry the positioning, and the categorical-semantics thread is one supporting strand.

\subsection{Notation diversity and representational bias}

Process mining uses many representational formalisms, Petri nets, process trees, causal nets,
directly-follows graphs, BPMN, Declare, and DCR graphs~\cite{deboisDeclarativeProcessMining2017},
each coupled to particular discovery algorithms. Three lines of work relate them. The first is
conversion-based translation with behavioural
preservation~\cite{kalenkovaProcessMiningUsing2017}. The second is van der Aalst's
representational-bias programme, which shows the target formalism constrains the discovery search
space and motivates sound-by-construction
classes~\cite{vanderaalstRepresentationalBiasProcess2011,%
  vanderaalstImprovingRepresentationalBias2012,%
  vanderaalstUsingFreeChoiceNets2021,%
vanderaalstFoundationsProcessDiscovery2022}. The third is the categorical-semantics tradition, in which Petri
nets present free symmetric monoidal categories~\cite{meseguerPetriNetsAre1990}, extended via open
Petri nets~\cite{baezOpenPetriNets2020} and string-diagrammatic
PROPs~\cite{bonchiDiagrammaticAlgebraLinear2019}. A common alternative is to route every notation through a single pivot such as workflow
nets~\cite{vanderaalstVerificationWorkflowNets1997}. That relocates the difficulty rather than
settling it, because the translations into the pivot are themselves uneven in rigour, they remain
incomplete on constructs such as the BPMN OR-join~\cite{dijkmanSemanticsAnalysisBusiness2008}, and
the comparison made afterwards is still a comparison of traces.
The present work adds a common compositional setting for the four formalisms it treats,
generalising conversion-based translation between BPMN, Petri nets, causal nets, and process
trees~\cite{kalenkovaProcessMiningUsing2017} to a single setting. These formalisms are
produced and consumed by process-mining tools such as ProM and PM4Py, so a shared framework also
serves the interchange between tools that emit different notations.

\subsection{Categorical and compositional process semantics}

Treating Petri nets as generators of symmetric monoidal categories originates with
\citet{meseguerPetriNetsAre1990}. Reading interface composition as an algebra of
connectors is the Montanari connector-algebra tradition, from the stateless connectors of
\citet{bruniBasicAlgebraStateless2006} through the concatenable-process
axiomatisation of \citet{sassoneAxiomatizationAlgebraPetri1996} and the classification of
models of concurrency of \citet{sassoneModelsConcurrencyClassification1996}.
Cospan composition here is the process-mining instance of that connector-composition tradition, with
the connectors read off event-log structure rather than fixed in advance. \citet{baezOpenPetriNets2020} extended
this to open Petri nets via decorated cospans. The present construction specialises that framework,
decorating cospan apices with process-mining data (activity labels, binding constraints) and
restricting to the set of runs $\mathbf{F}(\Sigma)$, the closed connected diagrams, to exclude non-process-like
wirings. \citet{bonchiDiagrammaticAlgebraLinear2019} unify nets and process calculi through
string-diagrammatic PROPs. Our aim is complementary, recovering partial-order structure from
event-log data rather than reasoning equationally about concurrency laws.
Similarly, \citet{lechenneCompositionalFrameworkPetri2024}
equip Petri nets with open-ended interfaces and a PROP-based graphical language in order to compute
reachability compositionally. Like the open-net line, the framework treats nets alone, with no
event-log or cross-notation discussion.

Outside the categorical approach, the Heraklit programme of \citet{fettkeSystemsMiningHeraklit2022} composes
Petri-net modules along left and right interfaces by fusing equally labelled boundary elements, an
associative composition proposed for process mining, and
they have recently shown that the partially ordered runs of a composed system are exactly the
composites of the runs of its parts~\cite{fettkeCompositionalitySystemsPartially2026}. This is the
compositionality that a hypergraph category provides by construction. Interface fusion is cospan
composition, and the run-level statement then holds uniformly across the four notations treated
here rather than being proved per formalism.

\citet{haymanUnfoldingGeneralPetri2008} showed that the cofree unfolding of a
general Petri net is recoverable only up to a symmetry on the folding morphism's kernel. The cospan
construction encounters the same multiplicity phenomenon, non-uniqueness of mediating morphisms from
indistinguishable tokens, but computes its quotient directly from generator boundaries, a form
estimable from event-log traces without solving a morphism-lifting problem.

\subsection{Partial orders, unfoldings, and untangling}

Partial-order-based process mining, spanning partially ordered event data, discovery,
and conformance, is surveyed by \citet{leemansPartialorderbasedProcessMining2023}.
\citet{Unfoldings2008} build finite complete prefixes of occurrence-net
unfoldings for model checking. Like the present work they take partial orders as primary, but
address the forward problem (verifying a given net) and need boundedness for finiteness, whereas the
cospan construction works directly from generators, essential where object-centric place
multiplicities are unbounded.

\citet{polyvyanyyUntanglingsNovelApproach2015} extract untanglings, representative
replay-compatible executions of a marked net. In the cospan setting a single string diagram
represents the whole concurrency class of executions with the same causal structure, and building
all diagrams from a signature enumerates all untanglings, most cleanly for 1-safe nets.
The three approaches differ in how much they enumerate. Unfoldings are fully enumerative
(a node per configuration), untanglings keep a pruned representative family, and the cospan
signature is structural, encoding causal roles rather than runs, with size
$\propto |A|\times|\mathrm{Contexts}|$ independent of run count and specific partial orders produced
on demand by composition.
The same interleaving blow-up motivates conformance checking over partially ordered
traces. Partially ordered alignments were introduced for this purpose by
\citet{luConformanceCheckingBased2015}. \citet{leemansPartiallyOrderedStochastic2025} encode observed
behaviour as labelled partial orders precisely to avoid enumerating the interleavings of concurrent
activities. The diagrams here are a model-side counterpart, each representing its concurrency class
directly, so that comparison can stay at the partial-order level.

The same frequency-aware agenda extends to stochastic process discovery. \citet{leemansStochasticProcessDiscovery2024} study when this discovery can be done optimally, and
\citet{liDiscoveringStochasticCausal2025} discover stochastic causal nets, a stochastic
counterpart to the object-centric causal nets treated in Section~\ref{sec:causal-nets}. The
construction here is qualitative, and decorating its generators with execution weights is a natural
extension.

Model-to-model comparison over event structures has also been studied. \citet{armas-cervantesBehavioralComparisonProcess2014} compare process models through
canonically reduced event structures and report their behavioural differences. The certificate of
Section~\ref{sec:conversion-framework} plays the same role across notations, and its canonical form
is fixed by the construction rather than by a reduction procedure.

\subsection{Process trees and series-parallel structure}\label{subsec:pt-sp-related}

Process trees are a central block-structured notation, used both as a model class and as a discovery
intermediate~\cite{leemansDiscoveringBlockStructuredProcess2013,buijsQualityDimensionsProcess2014}.
Their sequence and parallel operators are the classical series and parallel compositions of partial
orders. Finite series-parallel (SP) posets, the $N$-free finite
posets~\cite{valdesRecognitionSeriesParallel1982,mohringComputationallyTractableClasses1989}, have a
finitely axiomatisable pomset theory~\cite{gischerEquationalTheoryPomsets1988} linking them to the
concurrency semantics of \citet{prattModelingConcurrencyPartial1986} and
\citet{lodayaSeriesParallelLanguages2000}. The process-mining literature invokes this
correspondence implicitly through block-structured syntax.
Appendix~\ref{sec:pt-extras} states it as an explicit occurrence-level representation theorem
(Corollary~\ref{cor:pt-sp-iff}), with Theorem~\ref{thm:pt-as-sp-language} separating the $\to,+$
fragment from the full language, where XOR contributes set-union and loop bounded unfoldings.
The series-parallel restriction has also been relaxed at the notation level. POWL
models~\cite{kouraniPOWLPartiallyOrdered2023} compose submodels along arbitrary partial orders
rather than series and parallel operators alone, remain sound by construction as a subclass of
workflow nets, and have been extended to non-block-structured
choice~\cite{kouraniUnlockingNonBlockStructuredDecisions2025}. A language-preserving translation from safe
and sound workflow nets into POWL exists~\cite{kouraniTranslatingWorkflowNets2025}. POWL thus
removes at the syntax level the $N$-freeness that Corollary~\ref{cor:pt-sp-iff} characterises at
the semantic level, and its translation guarantee is stated on trace languages, whereas the
certificate of Section~\ref{sec:conversion-framework} compares structure, which by
Lemma~\ref{lem:equiv-strictness} is strictly finer.

\subsection{BPMN and formal process semantics}\label{subsec:bpmn-related}

Formal BPMN semantics are mainly operational. \citet{dijkmanSemanticsAnalysisBusiness2008} normalise a model and encode it into a Petri net
whose reachable markings are its valid token states. Section~\ref{sec:bpmn} uses no soundness notion. The soundness taxonomy of the workflow-net
literature~\cite{vanderaalstSoundnessWorkflowNets2011} has decision problems that are
\textsc{ExpSpace}- and \textsc{PSpace}-complete~\cite{blondinComplexitySoundnessWorkflow2022}. The
cospan construction is denotational instead. It extracts the interface structure with no firing
strategy or structural precondition, handling OR gateways by finite XOR/AND explosion, so it
avoids those decision problems, which then bound only the composition it defers.

Object-centric BPMN has been approached differently before. \citet{seidelObjectcentricBPMNProcess2024} give a fragment-based dialect with data objects and object
lifecycles and informal semantics, within the object-centric process-mining
programme~\cite{vanderaalstObjectCentricProcessMining2023,bertiAdvancementsChallengesObjectCentric2023}. Adjacent formal
extensions target the artifact-centric~\cite{lohmannArtifactCentricModelingUsing2012} and
data-aware~\cite{calvaneseFormalModelingSMTBased2019} paradigms. None types BPMN's control flow with object
types, as Section~\ref{sec:bpmn} does.

\subsection{Object-centric process mining}

Object-centric process mining sharpens the inter-notation problem.
OCEL~2.0~\cite{bertiOCELObjectCentricEvent2024} dropped the single-case constraint and object-centric Petri
nets~\cite{vanderaalstDiscoveringObjectCentricPetri2020} added typed tokens, with more
expensive reachability and less mature conformance. The standard does not natively carry
hierarchical structure, temporal inter-object relations, or explicit
state~\cite{khayatbashiAdvancingObjectcentricProcess2026,bertiStateAwareObjectCentricProcess2025}.
\citet{lissObjectCentricCausalNets2025} extend causal nets to typed binding sets, which
the cospan construction accommodates directly (Section~\ref{sec:causal-nets}). In the cospan
framework object types are wire labels. Boundary-matching composability enforces type-compatibility
structurally, the categorical analogue of a typed slice, with no reachability bookkeeping, and the
same mechanism covers the object-centric causal nets and process trees of Liss et al.\ and
\citet{vandettenDiscoveringCompactLive2024} unchanged. These object-centric
concerns are active in the modelling community as well as in mining. Data-aware enterprise process
modelling in MERODE~\cite{snoeckSupportingDataawareProcesses2023} and the object-centric
event-logging and data-centric modelling terminology of
iDOCEM~\cite{verbruggenIDOCEMDefiningCommon2024} formalise objects and their life-cycles at the
conceptual level, and model-driven management of BPMN-based process
families~\cite{delgadoModeldrivenManagementBPMNbased2022} organises related models within one
notation. The framework developed here is complementary and extends these by comparing and
translating models across the four notations at once, with object types carried structurally on the
wires rather than reconstructed from a conceptual schema.
